Wu et al.\ showed that a linear differential equation can be written as the ground state of a physics-informed effective Hamiltonian, assembled from circuits whose gate count grows polynomially in the number of qubits \cite{wu2025pihm}. That gate count prices one query to the Hamiltonian. The number of queries is the degree of the spectral filter, which grows as $\sqrt{\alpha_H/\Delta}\,\log(1/\varepsilon)$ for an encoding of subnormalisation $\alpha_H$, a spectral gap $\Delta$ and a target suppression $\varepsilon$, and the two problems that Wu et al.'s Discussion left open, the gap's dependence on the resolution and the cost of differentiation, are its denominator and, on our reading, its numerator. We asked how the number of queries grows with the resolution and how much of that growth is the encoder's, whether it can be lowered by changing the residual rather than the circuit, and whether the answer survives the step to nonlinear equations.

Separating the encoder's cost from the Hamiltonian's required a fixed construction, the standard form, in which the state holds the field alone and every derivative is eliminated through the dense Chebyshev differentiation matrix $\G$, with the printed defects removed. With Wu et al.'s printed choices restored it reproduces their eleven \gls{ode} examples. An exact factorisation of $\G$ lowers the estimated degree on the critically damped oscillator R2a at $n=8$ from $\res[2,sci]{stateprep/R2a/standard/n8/T3/filter=minimax/encoder=published}{degree}$ to $\res[2,sci]{stateprep/R2a/standard/n8/T3/filter=minimax}{degree}$, which is the encoder's share. The remainder belongs to the Hamiltonian. Each derivative eliminated through $\G$ multiplies the residual's norm by a factor of order $N^2$, for $N=2^n$ Chebyshev coefficients, so $\|\Hstd\|=\Theta(N^{4k})$ for an order-$k$ equation with constant coefficients, and the degree grows as $N^{2k}$ under the exact encoding and at the ideal-encoding bound $\alpha_R=\|R\|_2$ alike. The ideal bound contains no circuit, so this growth is the cost of differentiation itself, paid in queries.

Lowering that cost means changing the residual while keeping its kernel. The descriptor form keeps $\G$ as the quotient of two banded factors, $\G=\Bt^{-1}\Dt$, and carries the intermediate derivatives as unknowns tied by them, which makes the residual banded and leaves its kernel unchanged. Banded differentiation is the ultraspherical method's \cite{olver2013}, and carrying derivatives so as not to square a condition number is the idea of first-order-system least squares \cite{cai1994fosls}; what we add is their transfer to a block-encoded Hamiltonian and an account of its cost. The Hamiltonian's norm falls to $\Theta(N^2)$, and the residual compiles to a single \gls{lcu} whose ancilla grow as $\log n$. On a regular problem, one axis with polynomial coefficients, a leading coefficient without zeros on $[-1,1]$, analytic data and a solution unique up to scale, the gap is proved bounded below independently of $N$ and the degree grows as about $N$. On R2a at $n=8$ it is $\res{stateprep/R2a/descriptor/n8/T3/filter=minimax}{degree}$, a figure exact emulation confirms, and $\res{stateprep/R2a/descriptor/n8/T3/filter=minimax}{ideal_degree}$ at the ideal bound, where the standard form needs $\res[2,sci]{stateprep/R2a/standard/n8/T3/filter=minimax}{degree}$ and $\res[2,sci]{stateprep/R2a/standard/n8/T3/filter=minimax}{ideal_degree}$ (resource estimates).

The price of carrying the derivatives is paid in the gap. Their norms leave the Hamiltonian and enter the state's own norm, so the gap closes wherever the equation fails to control a derivative the state carries, as at a zero of the leading coefficient, the mark of a singular problem, and on a second axis. On one axis an a priori block rescaling restores the gap without ancilla, and the degree on singular problems then grows as about $N^2$, although on Legendre's equation the standard form needs the smaller degree at the ideal bound. On two axes the rescaled gap still falls: on the heat equation the descriptor's degree grows as $N^{\res[2]{stateprep/R5b/descriptor/growth/T3/filter=minimax}{degree}}$ per axis between four and five qubits per axis, and at four the standard form's is the lower, as built and at the ideal bound. On constant-coefficient problems a descriptor query also costs more two-qubit gates than a standard one at the sizes built and verified, so the advantage there lies in the number of queries.

Nonlinear equations needed a unique ground state before they could be prepared at all. Wu et al.'s doubled space encodes a quadratic term exactly, but its kernel has dimension $\res{build/R6a/standard/n5}{kernel_count}$ in a space of $\res{build/R6a/standard/n5}{dimension}$ at $n=5$, so a filter returns a projection onto it. A Carleman lift is instead a coupled linear system, carried like any other with several fields on one state, and its kernel is one-dimensional at every order run. The lift is prepared in both forms up to two-dimensional incompressible Navier--Stokes on $12$ to $48$ Fourier modes, by exact emulation and resource estimate, the descriptor's time axis lowering the degree from four time qubits upwards. A trusted kernel does not, however, certify the lift. On Burgers' equation at strength $a/\nu=10$, and on Navier--Stokes at $a/\nu=30$ from $28$ modes, the lift diverges while its kernel stays trusted, so the filter faithfully prepares a linear problem that is no longer the nonlinear one. Where the lift converges, its error follows $a/\nu$ rather than the worst-case ratio $\rho$ of the convergence theorem, whose condition $\rho<1$ is sufficient but not necessary.

All of this is classical computation, and nothing ran on hardware; every circuit is verified against the operator assembled from the problem, and where circuits were simulated they match exact emulation within $10^{-10}$. Neither form is polylogarithmic in $N$. On a regular problem a preparation from an initial state of overlap $\gamma$ with the solution costs $\tilde\Theta(N\,\mathrm{poly}(n)/\gamma)$ gates in the descriptor form against $\tilde\Theta(N^{2k}\,\mathrm{poly}(n)/\gamma)$ in the standard form, while on one axis a classical spectral solver needs only $\mathcal{O}(N)$ operations.

\label{sec:future-work}
Any advantage this route offers must therefore come on several axes, where the descriptor's gap falls, and it needs a pipeline that no longer takes its spectral edge from a classical solution of the problem. The gap is the more pressing of these limits, because it decides whether the reformulation can compete at all. Its fall is consistent with a loss of coercivity, an inference rather than a measurement: the datum rows weigh a boundary function equally at every frequency, too weakly to control the high-frequency states that the equation alone leaves unpenalised. First-order-system least squares answers the same failure with boundary residuals in fractional-order norms, which suggests weighting the datum rows of Laplace's equation by a diagonal that grows with the Chebyshev index. Such weights leave an exact kernel unchanged, so measuring the rescaled gap, and the norm the weights add, at three to five qubits per axis is a cheap first test. The alternative is the full ultraspherical reduction, a banded residual on the field alone with no block register, which was neither built nor compared here; one problem of each class, R2a, Legendre's equation and the heat equation, would show whether it keeps a norm of $\Theta(N^2)$ without eroding the gap.

The spectral edge is the other input still borrowed from the answer. An underestimated gap raises the degree only as the inverse square root of the bound, so a lower bound of the right order in $N$ would preserve every law measured here. On regular problems that order is proved by compactness, which yields no constant; a quantitative version for constant coefficients, replacing the compactness step with the explicit solution operator of the equation and its constraints and checked against the gaps measured on R1 and R2, is a well-posed first step. Singular and two-axis problems have no theorem to sharpen, and there the edge would have to be estimated on the circuit from the eigenphases of the qubitised walk, a harder task, since an initial state chosen without the answer has no reason to put weight near the edge.

The extensions leave narrower questions. The divergent lifts were caught only against the Galerkin solution, whereas comparing the fields decoded at consecutive Carleman orders needs no reference; whether that comparison separates convergence from divergence can be checked classically on the references already computed, which would also locate Burgers' boundary between $a/\nu=3$ and $10$. Native complex preparation needs complex loaders rather than a new filter, since a complex residual still gives a Hermitian $R^\dagger R$, and $z'=-\mathrm{i}\pi z$, realified as the oscillator R7, is the natural first comparison. The data-informed Hamiltonians of the research proposal fit the same residual, an observation $f(x_i)=y_i$ entering as the row $y_sf(x_i)-y_if(x_s)=0$ against the regular datum $f(x_s)=y_s$, though noisy observations replace the kernel with a least-squares fit whose gap has not been studied. The costs above stop at the prepared state, and a shot-count study of Wu et al.'s interferometric readout on R2a would supply the factor they omit. Hardware is further away than the circuit widths suggest: the smallest preparation, the decay $f'+2f=0$ (R1) at $n=2$ in the descriptor form, a circuit of only $\res{stateprep/R1/descriptor/n2/T1/filter=minimax}{n_qubits}$ qubits, applies $\res{stateprep/R1/descriptor/n2/T1/filter=minimax}{degree}$ queries of $\res{encode/R1/descriptor/n2}{cx_count}$ two-qubit gates each, which would place a first demonstration on early error-corrected machines.

The conclusion that carries beyond this method concerns where a ground-state solver's cost is set. The degree of a polynomial filter grows at least as the square root of the encoded operator's norm over its gap, so once the encoding is tight only the operator can be improved, and a reformulation that keeps the kernel may change both its norm and its gap. Here that meant conditioning the residual before encoding it. In the descriptor form the conditioning of differentiation moves out of the Hamiltonian's norm and into its gap, so the reformulation gains wherever the equation controls every derivative it carries, which a rescaling secures on one axis and nothing yet secures on two.
